\section{Workloads and Experimental Design}\label{sec:workloads}
\subsection{Sources}
\label{sec:sources}

LSSP's primary benchmark corpus draws from three independently released
workload sources: BurstGPT \citep{burstgpt2025}, Microsoft's official Azure
2024 LLM-inference trace \citep{azure2024trace,stojkovic2025dynamollm}, and
Alibaba Cloud's Bailian/Qwen anonymized usage-trace
release \citep{bailianqwen2026}. Azure's 2023 conversation/code traces
\citep{patel2024splitwise} and windows from an earlier discriminability
pilot (\S\ref{sec:stage0}) contribute additional context but are not part
of the primary 120-window matrix beyond the 10-per-source pilot windows
reused as described below. TraceLab \citep{tracelab2026} is retained as an
out-of-distribution, coding-agent workload family for future expansion,
not as one of the three primary sources, because its independence from a
prior 512-window public sweep cannot be fully verified from public
artifacts alone.

BurstGPT's own reported
collection methodology logs traffic from a regional customer of Microsoft's
Azure OpenAI GPT service, and this project's Azure trace is Microsoft's own
official release of internal Azure LLM-inference-service traffic. Both
therefore sit on the same underlying cloud provider's infrastructure. We
describe BurstGPT and Azure as distinct, independently released workload
sources -- different collection methodology, different customer
population, different release provenance and version history -- and we do
not describe them as independent providers. Bailian/Qwen (Alibaba Cloud)
is a genuinely different underlying platform from both.

\subsection{Representation and Provenance}
\label{sec:provenance}

Each source is ingested through a dedicated, field-level-audited adapter that maps the source's native
schema onto a common internal request representation (arrival time, input
tokens, output tokens, total tokens, session identifier where available,
model class). Native-field provenance is separate from benchmark synthesis: absent tenant, KV-reuse and failure fields are not inferred, whereas the benchmark supplies priority/SLO parameters for policies and metrics where necessary. These are synthetic experimental attributes, not source-observed service guarantees. License status for each source was verified before any confirmatory use of a real, full-scale sample (for example, BurstGPT is under CC-BY-4.0).

\subsection{Window Construction}
\label{sec:windows}

The benchmark corpus freezes exactly 120 windows of 200 requests each: 40
per source (\texttt{azure\_llm\_2024}, \texttt{bailian\_qwen},
\texttt{burstgpt}), each composed of 10 windows reused verbatim from the
discriminability pilot (\S\ref{sec:stage0}) plus 30 new windows drawn by a
deterministic, outcome-blind, source-symmetric pinned-extension sampler
that uses only valid-row index ranges and fixed seeds -- never scheduler
metrics, policy identities, or the pilot's tie/non-tie outcomes (verified
by structural audit to contain no BurstGPT-outcome-targeting signal).
Windows are assigned to EARLY/MIDDLE/LATE chronology strata by
deterministic relative chronology (60/30/30 windows respectively); the
reused pilot windows are all EARLY. The freeze was independently
validated -- unique window identities, exact per-source and reused/new
counts, fixed \texttt{request\_count = 200}, no prohibited overlap, valid
chronology, and absence of any scheduler/policy/load-calibration outcome
field in the selection rule -- and recorded with a canonical content hash
(excluding generation timestamp; full values and verification records are
in the preserved artifact hash ledger) before any load calibration or
scheduler outcome was generated.

\subsection{Descriptors}
\label{sec:descriptors}

Each window is described by a 28-feature descriptor vector spanning prompt-
and output-token-length statistics (mean, p90, long-prompt fraction at the
512/2048/8192-token thresholds), arrival-timing statistics (mean arrival
rate, inter-arrival percentiles), and derived load-pressure proxies. Three
deterministically redundant pressure-proxy features ($R^2=1.0$ against
other included features) are retained in the underlying dataset but
excluded from the classifier input used in \S\ref{sec:separability}.
Prompt-token figures compose differently across sources for architectural
reasons: TraceLab's and Azure's prompt-token fields are potentially
multi-turn-cumulative context, while BurstGPT's is a single per-request
figure. This is a property of how each source's underlying workload is
architected (a growing agent/conversation session versus a single stateless
API call), not an adapter artifact, and it is preserved as an explicit
caveat on every prompt-length-based claim below.


\subsection{Characterization scope}\label{sec:cross-source-diffs}\label{sec:separability}\label{sec:drivers}
A separate descriptor study used 400 windows across four classes, including
TraceLab. It is not part of the 120-window outcome matrix. Its classifier
results and feature-attribution claims require their own provenance and
are not evidence that a particular feature causes ranking reversals. This
draft retains the source and prompt-semantics caveats without using those
classifier accuracies as central portability evidence.
\subsection{Operating-Region Calibration}
\label{sec:calibration}

A FIFO-only calibration pass establishes six operating regions --
\texttt{LOW} (load factor 0.5), \texttt{PRE\_KNEE} (0.8), \texttt{KNEE}
(1.0), \texttt{POST\_KNEE} (1.1), \texttt{OVERLOAD} (1.2), and
\texttt{HIGH\_PRESSURE} (1.5) -- across all 120 frozen windows (720
calibration cells total; 720 valid, deterministic, 0 failed/missing/
duplicated). We report this as calibration provenance, not a comparative
scheduler result: FIFO is the calibration reference policy, and the
calibration's role is to fix load-factor-to-region mappings before any
other policy is executed against them, not to characterize FIFO's relative
performance. The preserved 720-cell raw FIFO calibration manifest contains no zero-completion cells. Earlier manuscript prose reporting 301 such records is inconsistent with this artifact and is not retained as a finding. The distinction between calibration and campaign configuration is preserved: a calibration record is not an outcome for every policy under the final campaign.

\subsection{Paired Design and Repetitions}
\label{sec:paired}

Because the simulator is fully deterministic (\S\ref{sec:execution-model}),
a single execution per (window, policy, region) cell is sufficient: each of
the $9{,}360$ configurations yields one scientific outcome, and its rep0/rep1
pair ($18{,}720$ executions in total) is a deterministic verification pass
confirming bit-identical output, not an independent stochastic replicate.
Ranking-uncertainty quantification (\S\ref{sec:results}) instead comes from
resampling over the workload windows themselves (block
bootstrap), not from repeated stochastic runs of the same cell.


\subsection{Protocol freeze, pilot and overlap}\label{sec:prereg}\label{sec:stage0}\label{sec:evidence-boundaries}\label{sec:overview}
The repository preserves the workload, panel, calibration, campaign and
analysis freezes and their hashes. The earlier discriminability pilot
motivated the larger panel; it is not confirmatory ranking evidence.
The principal outcome matrix is fixed. Subsequent metric and SLO analyses
are labeled extensions, with their original configurations unchanged.
Reused pilot windows are disclosed rather than counted as an independent
replication of the pilot. Historical source and freeze identities are
retained even where later prose or engineering audits qualify their meaning.

A related load-scaling manuscript and a separate module-intervention
project share research lineage. This paper does not claim their results
or intervention-based attribution as its own. The earlier 60-window
project's exact byte-level overlap is not fully reconstructed; final
overlap review remains necessary. The present sampler is distinct, but
that alone is not a proof of zero overlap.

\subsection{Window and trace eligibility}\label{sec:window-size}
Every principal window contains 200 requests; a sweep of requests per
window has not been performed. The recovery analysis varies the number
of windows, not requests per window. Source-reported token counts are not
retokenized model inputs. All 40 Azure and 40 Bailian window medians exceed
256 source tokens. The later whole-prefill engineering contract, bounded
at 256 prompt tokens and 16 active sequences, therefore does not cover
these traces. No requests were truncated or filtered for this manuscript
preparation, and no mapping to a modern model/tokenizer context limit is
silently assumed.

\subsection{Temporal partitions}\label{sec:temporal-design}
The manifest assigns 20 EARLY, 10 MIDDLE and 10 LATE windows per source
(60/30/30 over all three sources). Thus 60/30/30 is not a BurstGPT-only
count. BurstGPT uses relative-chronology terciles and an additional bisect
sensitivity split. Azure uses the frozen 2024-05-15 calendar boundary;
Bailian is labeled relative chronology only. These definitions remain
unchanged and are interpreted separately.
\subsection{Recorded prompt-length coverage}
\begin{table}[t]\centering
\caption{Range of window-level median source-reported prompts in the
frozen 40-window/source corpus. These are not common-tokenizer counts.}
\label{tab:prompt-coverage}
\begin{tabular}{lrrr}\toprule
Source & Minimum median & Maximum median & Medians above 256\\\midrule
azure\_llm\_2024 & 628.5 & 1595 & 40/40\\
bailian\_qwen & 343.5 & 1150 & 40/40\\
burstgpt & 140.5 & 620.5 & 9/40\\
\bottomrule\end{tabular}\end{table}
These summaries demonstrate a coverage gap for bounded short-prompt
serving checks, not an exact request-exclusion fraction. A window median
does not identify every request above a threshold, and source length
fields do not recover absent prompt text. No benchmark text is synthesized
here to claim tokenizer equivalence. Model, tokenizer and context eligibility
for prospective native traces require separate decisions and measurements.
